\section{Schwarzschild slices: original charges and sharpness}
\label{q:converse:sec}

We now prove the converse in the strong-decay class. Start with a spacelike
hypersurface in a Schwarzschild--Tangherlini exterior of mass $M>0$. The geometry of
its horizon determines the area, but the mass in the Penrose inequality
is computed in the hypersurface's own asymptotic coordinates.  We must
therefore identify its original ADM charges, including when the slice
is asymptotically boosted.

The proof has three steps.  The induced-data decay forces a single
spacelike plane at infinity.  We compute that plane's fluxes and show
that changing its height and its spatial coordinates contributes zero
in the limit.  Finally, the horizon area gives equality whenever the
hypersurface has the full-cut minimality required in the rigidity
class.  Radial compactly supported changes of the static slice will
then supply non-time-symmetric examples and verify all those cut and
outermostness hypotheses directly.

Throughout this section $n\ge3$, $k=n-1$, $p=n-2=k-1$,
$\omega=|S^{n-1}|$, and $r_h=(2M)^{1/p}$.
We use $\mathfrak r$ for the areal radius and $s=|y|$
for the isotropic radius.  In static isotropic coordinates $(b,y)$ the
spacetime metric is
\begin{equation}
 \mathbf G=-\alpha(s)^2\,db^2+h_0,
 \qquad
 h_0=\left(1+\frac{M}{2s^p}\right)^{4/p}\delta,
 \qquad
 \alpha(s)=\frac{1-M/(2s^p)}{1+M/(2s^p)}.
 \label{q:converse:eq:static-isotropic}
\end{equation}
We retain the future-normal convention
$K(U,V)=\mathbf G(\boldsymbol\nabla_U\mathbf n,V)$, and the energy
and momentum factors $1/(2k\omega)$ and $1/(k\omega)$, respectively.
The static Killing field $\xi=\partial_b$ is future timelike in the open
exterior.  The isotropic spatial coordinates extend continuously to its
horizon by the angular coordinates and the limiting isotropic radius.

We use the compact-complement convention of Definition~\ref{def:data}:
outside a compact subset of $\Omega$, there is just the given coordinate
end $\{|x|>R_0\}$.  For an embedding of this end, ``approaches spatial
infinity'' means
\begin{equation}
 |x_j|\longrightarrow\infty
 \quad\Longrightarrow\quad
 |y(x_j)|\longrightarrow\infty
 \quad\hbox{for every sequence in the end.}
 \label{q:converse:eq:spatial-infinity}
\end{equation}
This is the proper escape condition used below.  An assertion about the
existence of a single escaping sequence would not suffice.

\subsection{The asymptotic plane is forced by the induced data}

\begin{lemma}[An arbitrary admissible slice has a plane at infinity]
\label{q:converse:lem:plane}
Let a smooth spacelike embedding of $\Omega$ have its interior in the
open exterior and extend smoothly to its compact boundary in the
regular horizon extension.  Suppose its induced data obey
\[
 g_{ij}-\delta_{ij}=O_6(r^{-q}),\qquad
 K_{ij}=O_5(r^{-1-q}),\qquad
 \frac p2<q<p,
 \qquad r=|x|,
\]
with the future-normal convention stated above.
Assume the compact-complement condition and
Equation~\eqref{q:converse:eq:spatial-infinity}.
For every $q_0$ with $p/2<q_0<q$, the static Killing field decomposes as
\begin{equation}
 \xi=u\mathbf n+X,\qquad
 (u,X)=(u_\infty,X_\infty)+O_2(r^{-q_0}),\qquad
 u_\infty>0,\qquad u_\infty^2-|X_\infty|^2=1.
 \label{q:converse:eq:kid-limit}
\end{equation}
After a constant time-orientation-preserving Lorentz change of ambient
coordinates, the embedded end is a single graph $T=f(z)$ over a full
distant spatial coordinate tail.  After a constant orthogonal change of
the $x$ coordinates, its coordinate change and height obey
\begin{equation}
 \begin{aligned}
 z(x)&=x+\psi(x),&\quad \psi=o(r),&\quad
 D\psi=O_1(r^{-q_0}),\\
 f(z)&=o(|z|),& Df&=O_1(|z|^{-q_0}).
 \end{aligned}
 \label{q:converse:eq:graph-estimates}
\end{equation}
\end{lemma}

\begin{proof}
We first derive estimates in the original $x$ coordinates, without using
the ambient spatial projection as a global coordinate system.  Since
$\xi$ and $\mathbf n$ are future timelike,
\[
 u=-\mathbf G(\xi,\mathbf n)>0,
 \qquad W:=u^2-|X|_g^2=\alpha(s)^2>0.
\]
Killing translation preserves the induced metric and second form.
Its tangential and normal parts therefore give the vacuum Killing
initial-data equations
\begin{align}
 \nabla_iX_j+\nabla_jX_i&=-2uK_{ij},
 \label{q:converse:eq:kid-first}\\
 \nabla_i\nabla_ju
 &=u\bigl(\Ric_{ij}+\tau K_{ij}-2K_i{}^\ell K_{\ell j}\bigr)
       -(\mathcal L_XK)_{ij}.
 \label{q:converse:eq:kid-second}
\end{align}
The sign of the first equation follows directly by restricting
$\mathcal L_\xi\mathbf G=0$ to two tangential vectors and using
$K(U,V)=\mathbf G(\boldsymbol\nabla_U\mathbf n,V)$.
For the second equation, the normal variation of $K$ with speed $u$ in
vacuum is
$\nabla^2u-u(\Ric+\tau K-2K^2)$ with this sign convention;
adding the tangential variation $\mathcal L_XK$ gives zero.
In particular, no stationary development has to be constructed in this
argument: the ambient vacuum metric is already given.

To see the precise differential estimate supplied by these equations,
expand
\[
 (\mathcal L_XK)_{ij}
 =X^\ell\nabla_\ell K_{ij}
   +K_{\ell j}\nabla_iX^\ell+K_{i\ell}\nabla_jX^\ell.
\]
Differentiating Equation~\eqref{q:converse:eq:kid-first}, adding the two
permutations with first derivative indices $i,j$, and subtracting the
third permutation expresses $\nabla_i\nabla_jX_\ell$ as curvature times
$X$ and the three derivatives
$-\nabla_i(uK_{j\ell})-\nabla_j(uK_{i\ell})+
\nabla_\ell(uK_{ij})$.  Commuting covariant derivatives accounts for
the curvature terms.  Thus both Hessians involve only first derivatives
of $u,X$, with coefficients $K$, and their values, with coefficients
among $\Ric$, the full curvature, $\nabla K$, and $K^2$.
Passing to coordinate derivatives adds coefficients $\Gamma(g)$ to
first derivatives and $D\Gamma(g)+\Gamma(g)^2$ to values.  The assumed
decay consequently gives, for $Y=(u,X)$,
\begin{equation}
 |D^2Y|\le C r^{-1-q}|DY|+C r^{-2-q}|Y|.
 \label{q:converse:eq:jet-bound}
\end{equation}

Applying Lemma~\ref{lem:static:end-jets} to
Equation~\eqref{q:converse:eq:jet-bound} gives
\[
 (u,X)=(u_\infty,X_\infty)+O_2(r^{-q_0}).
\]
Its hypothesis $q>1/2$ includes the slow decay permitted when $n=3$.
Equation~\eqref{q:converse:eq:spatial-infinity} gives $W\to1$, so taking
limits in $W=u^2-|X|_g^2$ proves
Equation~\eqref{q:converse:eq:kid-limit}.

We next justify the global coordinate assertion on a tail.
The spatial projection $\pi$ of the slice to $y$ is a local
diffeomorphism: a nonzero vector in its differential kernel would be
vertical and therefore timelike, whereas every nonzero tangent vector
of the slice is spacelike.  Choose a target radius $s_0$ so large that
the compact complement of the original coordinate tail, including its
inner coordinate sphere, projects inside $\{|y|<s_0\}$.
The map
\[
 \pi:\pi^{-1}(\{|y|>s_0\})\longrightarrow\{|y|>s_0\}
\]
is proper.  Indeed, the preimage of a compact target set cannot escape
through the intrinsic end by
Equation~\eqref{q:converse:eq:spatial-infinity}, and cannot approach the
excluded compact core by the choice of $s_0$.
Its image is open by the local inverse theorem and closed by
properness.  It is nonempty and the target is connected, so the image
is the whole tail.  A proper local diffeomorphism is a covering: its
fiber over a point is compact and discrete, hence finite; finitely many
local inverse neighborhoods can be shrunk simultaneously, and
properness excludes any additional preimages over that neighborhood.

Each component of this covering is one sheet, since
$\R^n\setminus\overline B_{s_0}$ is simply connected for $n\ge3$.
There is only one such component.  To check this last point, the
complement of the preimage tail is compact by the same escape
condition.  Each sheet is noncompact, and points on it with
$|y|\to\infty$ leave every intrinsic compact set.  Two sheets would
therefore give two distinct ends after removing that compact
complement.  This contradicts the stipulated single end.  We have
proved that $y$ is a global coordinate system on this tail, and that
the static time $b$ there is a single-valued function of $y$.

Let $\overline g=\pi^*h_0$.  Restricting $\xi^\flat$ to the slice gives
$X_g^\flat=-W\,db$, whence
\begin{equation}
 db=-\frac{X_g^\flat}{W},\qquad
 \overline g=g+\frac{X_g^\flat\otimes X_g^\flat}{W}.
 \label{q:converse:eq:base-pullback}
\end{equation}
By Equation~\eqref{q:converse:eq:kid-limit}, $\overline g$ tends to the
positive definite constant matrix
$\delta+X_\infty^\flat\otimes X_\infty^\flat$, with
$O_2(r^{-q_0})$ error.  Since $h_0$ tends to $\delta$, the identity
$\overline g=(D_xy)^T h_0(y)D_xy$ bounds both $D_xy$ and its inverse.
Integrating $D_xy$ along intrinsic radial paths yields $|y|\le Cr$.
Conversely, integrate the inverse derivative along target radial paths
starting at a fixed target sphere.  Its preimage is compact, giving
$r\le C|y|$.  Thus the two radii are comparable.

The connection transformation law is now available in genuine
coordinate charts:
\begin{equation}
 \partial_i\partial_jy^a
 =\Gamma(\overline g)^\ell_{ij}\partial_\ell y^a
  -\Gamma(h_0)^a_{bc}(y)\partial_i y^b\partial_j y^c.
 \label{q:converse:eq:coordinate-connection}
\end{equation}
Here $\Gamma(\overline g)=O_1(r^{-1-q_0})$ and
$\Gamma(h_0)=O_1(s^{-1-p})$.  Because $q_0<p$, bounded first
derivatives and comparable radii give
$D_x^2y=O_1(r^{-1-q_0})$.
Radial integration, followed by comparison on coordinate spheres,
therefore gives a single invertible constant matrix $B_0$ with
\[
 D_xy=B_0+O_1(r^{-q_0}).
\]
Equation~\eqref{q:converse:eq:base-pullback} similarly gives
$db=\beta+O_1(r^{-q_0})$, where
$\beta=-X_\infty^\flat$.  Taking limits in the induced metric shows
\begin{equation}
 B_0^TB_0-\beta\otimes\beta=I.
 \label{q:converse:eq:limiting-plane}
\end{equation}
Thus the linear map $v\mapsto(\beta(v),B_0v)$ is an isometric
embedding of Euclidean $\R^n$ as a spacelike Minkowski hyperplane.
A time-orientation-preserving Lorentz transformation takes this
hyperplane to $T=0$.  In the resulting coordinates $(T,z)$,
\[
 D_xT=O_1(r^{-q_0}),\qquad D_xz=Q+O_1(r^{-q_0}),\qquad Q^TQ=I.
\]
Integration gives $T=o(r)$ and $z=Qx+o(r)$.  Choose a fixed large
intrinsic sphere $r=R_1$ beyond which these coordinates are defined
and $D_xz$ is invertible, and then choose
$Z_0>\max_{r=R_1}|z|$.  The map
\[
 z:\{r>R_1,\ |z(x)|>Z_0\}\longrightarrow\{|z|>Z_0\}
\]
is proper by $z=Qx+o(r)$ and the exclusion of that inner sphere.
The preceding covering argument makes it a single coordinate chart.
This use of boosted coordinates is confined to the tail; it requires
no extension of the static time to the compact horizon boundary.
Rotating $x$ by $Q$ and using the chain rule proves
Equation~\eqref{q:converse:eq:graph-estimates}.
\end{proof}

\subsection{The charges of the limiting plane}

We have found the asymptotic plane without choosing a preferred
foliation of the original hypersurface.  Before estimating its
remaining height, we compute the plane's energy and momentum directly.
This fixes both the normalization and the sign of the momentum.
In the next lemma $x$ denotes Euclidean coordinates on the reference
plane; the coordinates on the given slice will return in the following
subsection.

\begin{lemma}[Charges of a boosted Schwarzschild--Tangherlini end]
\label{q:converse:lem:boost}
In the Schwarzschild--Tangherlini spacetime of geometric mass $M>0$,
let $b,y$ be static time and isotropic spatial coordinates.  Make the
Lorentz change
\begin{equation}\label{q:converse:eq:boost-coordinates}
 b=\gamma(T+vx^1),\qquad y^1=\gamma(x^1+vT),\qquad
 y^a=x^a\ (2\le a\le n),\qquad \gamma=(1-v^2)^{-1/2},
\end{equation}
where $0\le v<1$.  The induced data on $T=0$, oriented by the
future normal, have
\begin{equation}\label{q:converse:eq:boost-charges}
 E=\gamma M,\qquad P_1=\gamma Mv,\qquad P_a=0\ (a>1).
\end{equation}
The metric and second fundamental form have the homogeneous leading
terms computed below, with remainders $O_d(r^{-2p})$ and
$O_d(r^{-2p-1})$, respectively, for every fixed derivative order $d$.
\end{lemma}

\begin{proof}
Expanding Equation~\eqref{q:converse:eq:static-isotropic}, the leading
perturbation of Minkowski space is
$2M|y|^{-p}(db^2+p^{-1}|dy|^2)$.  On $T=0$ put
\[
 \rho^2=\gamma^2(x^1)^2+\sum_{a=2}^n(x^a)^2,\qquad f=\rho^{-p},
 \qquad A=\gamma^2(v^2+p^{-1}),\quad B=p^{-1},\quad
 C=kp^{-1}\gamma^2v.
\]
The nonzero relevant leading components of the spacetime perturbation
$\mathsf H$ are
$\mathsf H_{11}=2MAf$, $\mathsf H_{aa}=2MBf$, and
$\mathsf H_{01}=2MCf$.  Moreover
\[
 \partial_1f=-p\gamma^2x^1\rho^{-n},\qquad
 \partial_af=-px^a\rho^{-n},\qquad
 \partial_Tf=-p\gamma^2vx^1\rho^{-n}.
\]
Thus the energy flux vector $Q_i=\partial_j\mathsf H_{ij}
-\partial_i\mathsf H_{jj}$ satisfies
\[
 Q_1=-2MkB\partial_1f=2Mk\gamma^2x^1\rho^{-n},\qquad
 Q_a=-2M(A+pB)\partial_af=2Mp(A+1)x^a\rho^{-n}.
\]
Since $p(A+1)=k\gamma^2$, these expressions give the pointwise identity
\begin{equation}\label{q:converse:eq:boost-energy-vector}
 Q_i=2Mk\gamma^2x_i\rho^{-n}.
\end{equation}

The chosen sign of the second fundamental form gives, to first order,
\[
 K_{ij}=\tfrac12\bigl(\partial_T\mathsf H_{ij}
                 -\partial_i\mathsf H_{0j}
                 -\partial_j\mathsf H_{0i}\bigr).
\]
Direct substitution gives
\begin{align*}
 K_{11}&=M\gamma^4v(2p+1-pv^2)x^1\rho^{-n},&
 K_{aa}&=-M\gamma^2vx^1\rho^{-n},\\
 K_{1a}&=Mk\gamma^2vx^a\rho^{-n},&
 K_{ab}&=0\quad(a\ne b).
\end{align*}
Their Euclidean trace is
$M\gamma^4v(p+v^2)x^1\rho^{-n}$.  Consequently, for the leading
Euclidean trace reversal $\pi_{ij}=K_{ij}-(\tr_\delta K)\delta_{ij}$,
\begin{equation}\label{q:converse:eq:boost-momentum-tensor}
 \pi_{1j}=Mk\gamma^2v x_j\rho^{-n},\qquad
 \pi_{a1}=Mk\gamma^2vx^a\rho^{-n},\qquad
 \pi_{ab}=-Mk\gamma^4vx^1\rho^{-n}\delta_{ab}.
\end{equation}
Using the exact normal and the exact metric in the trace reversal
changes these formulas by $O_d(r^{-2p-1})$.  The energy remainder
has first derivatives of this same order.  Their sphere integrals
are $O(r^{-p})$, and therefore vanish for every $n\ge3$.

The ellipsoid $\{|\operatorname{diag}(\gamma,1,\ldots,1)x|<1\}$
has volume $\omega/(n\gamma)$.  Polar integration therefore gives
\begin{equation}\label{q:converse:eq:ellipsoid-integral}
 \int_{S^{n-1}}|\operatorname{diag}(\gamma,1,\ldots,1)\theta|^{-n}
 \,d\theta=\frac{\omega}{\gamma}.
\end{equation}
Equations~\eqref{q:converse:eq:boost-energy-vector} and
\eqref{q:converse:eq:boost-momentum-tensor}, divided by $2k\omega$ and
$k\omega$, give $E=M\gamma$ and $P_1=M\gamma v$.
The integrand of a transverse momentum row is a constant times
$x^1x^a\rho^{-n}/r$, whose integral vanishes by reflection.
\end{proof}

The calculation is for each fixed $v<1$.  No estimate uniform as
$v$ approaches $1$ is used.

\subsection{Height and coordinate changes leave the charges unchanged}

The asymptotic plane does not by itself settle the charge calculation:
the graph Hessian can decay more slowly than the model second form.
We now separate its trace-reversed Hessian, whose flux vanishes
identically, from the errors whose sphere integrals tend to zero.


\begin{proposition}[ADM charges of an admissible Schwarzschild slice]
\label{q:converse:prop:adm}
Under the hypotheses of Lemma~\ref{q:converse:lem:plane}, the ADM limits
of the original induced data exist and satisfy
\begin{equation}
 E=M u_\infty,\qquad P_i=-M\delta_{ij}X_\infty^j,
 \qquad E>0,\qquad E^2-|P|^2=M^2.
 \label{q:converse:eq:charges}
\end{equation}
In particular the invariant mass is $M$.  No parity assumption is
needed.
\end{proposition}

\begin{proof}
Use the coordinates supplied by Lemma~\ref{q:converse:lem:plane} and
write $R=|z|$.  A constant Lorentz transformation of
Equation~\eqref{q:converse:eq:static-isotropic} gives, on every sufficiently
small fixed cone $|T|<\varepsilon R$,
\begin{equation}
 \mathbf G_{\alpha\beta}
 =\eta_{\alpha\beta}+\mathsf H_{\alpha\beta}
       +O_j(R^{-2p}),\qquad
 \mathsf H=2M|y(T,z)|^{-p}
       \left(db^2+\frac1p|dy|^2\right).
 \label{q:converse:eq:ambient-leading}
\end{equation}
Here $\eta=-dT^2+|dz|^2$ and the differentials on the right are the
constant linear Lorentz differentials.  The small cone is chosen so
that $|y(T,z)|$ and $R$ are comparable.  All coordinate derivatives of
any fixed order have their corresponding decay, and
$\mathsf H$ is homogeneous of degree $-p$.
Define the leading induced plane tensors by
\begin{equation}
 h_{ij}(z)=\mathsf H_{ij}(0,z),\qquad
 L_{ij}(z)=\frac12\left(
 \partial_T\mathsf H_{ij}
 -\partial_i\mathsf H_{0j}-\partial_j\mathsf H_{0i}
 \right)(0,z).
 \label{q:converse:eq:plane-tensors}
\end{equation}
The exact reference plane has metric
$\delta+h+O_2(R^{-2p})$ and second form
$L+O_1(R^{-2p-1})$.  Its charges were computed in
Lemma~\ref{q:converse:lem:boost}.  We will show that precisely the same
terms survive in the ADM fluxes of the given hypersurface.

First consider its graph $T=f(z)$ in these coordinates.
Its tangent vectors are $e_i=\partial_i+f_i\partial_T$, so direct
restriction of Equation~\eqref{q:converse:eq:ambient-leading} gives
\[
 g^{\rm gr}_{ij}
 =\delta_{ij}-f_if_j+\mathsf H_{ij}(f,z)
  +\mathsf H_{0j}(f,z)f_i+\mathsf H_{0i}(f,z)f_j
  +\mathsf H_{00}(f,z)f_if_j+O_1(R^{-2p}).
\]
In particular
\begin{equation}
 g^{\rm gr}_{ij}=\delta_{ij}+h_{ij}+e^{\rm gr}_{ij},
 \qquad
 |e^{\rm gr}|+R|De^{\rm gr}|
 \le C\left(
 |f|R^{-p-1}+R^{-2q_0}+R^{-p-q_0}+R^{-2p}
 \right).
 \label{q:converse:eq:graph-metric-error}
\end{equation}
For example, the evaluation difference
$\mathsf H_{ij}(f,z)-\mathsf H_{ij}(0,z)$ has size
$O(|f|R^{-p-1})$; its derivative also has the term
$O(|Df|R^{-p-1})$, included in the displayed bound.
Since $f=o(R)$ and $2q_0>p$, this proves
$e^{\rm gr}=o_1(R^{-p})$.

For completeness, the normal formula supplies the needed second-form
estimate with its sign.  In flat spacetime the future normal covector
of the graph is
\[
 \mathbf n^\flat
 =\frac{-dT+df}{\sqrt{1-|Df|^2}},
 \qquad
 K^{\rm flat}_{ij}
 =\frac{f_{ij}}{\sqrt{1-|Df|^2}}.
\]
In the perturbed metric, compute instead from
$K_{ij}=-\mathbf G(\mathbf n,\boldsymbol\nabla_{e_i}e_j)$.
The coefficient of $f_{ij}$ changes by
$O(|Df|^2+R^{-p})$.  The connection term at $T=0$, with zero slope,
is exactly $L$ to first order.  Evaluating it at $T=f$ costs
$O(|f|R^{-p-2})$, inserting the graph slopes costs
$O(R^{-p-1}|Df|)$, and quadratic metric corrections cost
$O(R^{-2p-1})$.  It follows that
\begin{equation}
 \begin{split}
 K^{\rm gr}_{ij}&=L_{ij}+f_{ij}+e^{K}_{ij},\\
 |e^K|&\le C\left(
 |f|R^{-p-2}+R^{-p-q_0-1}
       +R^{-3q_0-1}+R^{-2p-1}\right)
       =o(R^{-p-1}).
 \end{split}
 \label{q:converse:eq:graph-second-error}
\end{equation}
The $R^{-3q_0-1}$ term is the flat nonlinear error
$O(|Df|^2|D^2f|)$.  Thus even when $f_{ij}$ decays more slowly than
the plane second form, its only nonnegligible contribution is an
ordinary coordinate Hessian.

It remains to evaluate fluxes in the original coordinates.  Set
$z=x+\psi(x)$ as in Equation~\eqref{q:converse:eq:graph-estimates} and
$F_0(x)=f(z(x))$.  Pulling back the Euclidean metric gives the exact
expression
\[
 (D_xz)^TD_xz
 =I+D\psi+(D\psi)^T+(D\psi)^TD\psi.
\]
The last term is $O_1(r^{-2q_0})$.  Homogeneity and $\psi=o(r)$ give
\[
 (z^*h)_{ij}=h_{ij}(x)+o_1(r^{-p}),
 \qquad
 (z^*L)_{ij}=L_{ij}(x)+o(r^{-p-1}).
\]
Indeed, evaluation changes contribute $\psi\,Dh$ and $\psi\,DL$;
tensor pullback contributes $D\psi\,h$ and $D\psi\,L$.
Differentiating the metric pullback also produces
$D^2\psi\,h$, of size $O(r^{-p-q_0-1})$.
The chain rule gives a useful exact identity for the Hessian:
\begin{equation}
 z^*(D_z^2f)_{ij}
 =\partial_i\partial_jF_0
   -(\partial_af)(z(x))\,\partial_i\partial_j\psi^a.
 \label{q:converse:eq:hessian-chain}
\end{equation}
Its final term is $O(r^{-2q_0-1})=o(r^{-p-1})$.
Combining Equations~\eqref{q:converse:eq:graph-metric-error}--
\eqref{q:converse:eq:hessian-chain} therefore gives
\begin{align}
 g_{ij}
 &=\delta_{ij}+h_{ij}(x)
       +\partial_i\psi_j+\partial_j\psi_i+e_{ij},
 &e&=o_1(r^{-p}),
 \label{q:converse:eq:original-metric}\\
 K_{ij}
 &=L_{ij}(x)+\partial_i\partial_jF_0+\widetilde e_{ij},
 &\widetilde e&=o(r^{-p-1}).
 \label{q:converse:eq:original-second}
\end{align}
For the momentum trace reversal, the original metric satisfies
$g-\delta=O(r^{-q_0})$ and $K=O(r^{-q_0-1})$.  Consequently
\[
 (\tr_gK)g_{ij}-(\tr_\delta K)\delta_{ij}
 =O(r^{-2q_0-1})=o(r^{-p-1}).
\]
Thus
\begin{equation}
 \begin{split}
 K_{ij}-(\tr_gK)g_{ij}
 ={}&L_{ij}-(\tr_\delta L)\delta_{ij}\\
 &+\partial_i\partial_jF_0-(\Delta F_0)\delta_{ij}
       +o(r^{-p-1}).
 \end{split}
 \label{q:converse:eq:momentum-decomposition}
\end{equation}

Both apparent extra fluxes vanish exactly.  The energy flux vector
of the symmetric-gradient term in
Equation~\eqref{q:converse:eq:original-metric} is
\[
 Q_i(\psi)=\Delta\psi_i-\partial_i\divg\psi,
 \qquad \partial_iQ_i(\psi)=0.
\]
Choose a smooth cutoff that is one outside a fixed large sphere and
zero on a smaller ball, and use it to extend $\psi$ smoothly over
$\R^n$.  The same formula for the extension is divergence free on
the entire ball bounded by every sufficiently large coordinate
sphere.  The divergence theorem gives
\begin{equation}
 \int_{S_r}Q_i(\psi)\,\nu_\delta^i\,dA_\delta=0.
 \label{q:converse:eq:gauge-flux}
\end{equation}
Similarly, each row of the trace-reversed Hessian obeys
\[
 \partial_j\bigl(\partial_i\partial_jF_0
             -(\Delta F_0)\delta_{ij}\bigr)=0.
\]
Smoothly extending $F_0$ by the same procedure proves
\begin{equation}
 \int_{S_r}\bigl(\partial_i\partial_jF_0
             -(\Delta F_0)\delta_{ij}\bigr)
              \nu_\delta^j\,dA_\delta=0
 \quad\hbox{for every }i.
 \label{q:converse:eq:height-flux}
\end{equation}
These are identities for the extended potentials; no decay assumption
on an interior extension, and no subtraction of an unknown inner
boundary flux, is involved.

Finally, a metric error $o_1(r^{-p})$ contributes $o(1)$ to the
energy flux because $|S_r|=O(r^{p+1})$.  A second-form error
$o(r^{-p-1})$ contributes $o(1)$ to momentum.  More explicitly the
largest product errors above have flux size $O(r^{p-2q_0})$; the
remaining ambient products have size $O(r^{-q_0})$ or $O(r^{-p})$,
and the evaluation errors have size
$O(|f|/r+|\psi|/r)=o(1)$.
All tend to zero.  Equations~\eqref{q:converse:eq:gauge-flux} and
\eqref{q:converse:eq:height-flux} leave precisely the plane fluxes.

Lemma~\ref{q:converse:lem:boost} computes these, with the present second-form
convention, as $E=\gamma M$, $P_1=\gamma vM$, and $P_a=0$ for $a>1$
under
\[
 b=\gamma(T+vz^1),\qquad
 y^1=\gamma(z^1+vT),\qquad y^a=z^a.
\]
On that plane $\xi=\gamma\partial_T-\gamma v\partial_{z^1}$, so its
limiting lapse and shift are $(\gamma,-\gamma v,0,\ldots,0)$.
Constant spatial rotations give
$E=Mu_\infty$ and $P=-MX_\infty$ in the original orthonormal end
frame.  Equation~\eqref{q:converse:eq:kid-limit} proves the final
identity in Equation~\eqref{q:converse:eq:charges}.
\end{proof}

\subsection{The horizon area and the converse equality}

\begin{lemma}[Geometry of a full horizon section]
\label{q:converse:lem:horizon-area}
Every full smooth spacelike section of the future horizon, including
the bifurcation sphere, has induced area
\begin{equation}
 A=\omega r_h^k=\omega(2M)^{k/p}.
 \label{q:converse:eq:horizon-area}
\end{equation}
If such a section is the boundary of a smooth spacelike hypersurface
whose interior lies in the exterior, its outgoing future expansion
is zero with the future-normal convention used here.
\end{lemma}

\begin{proof}
In horizon-regular null coordinates, restriction of the spacetime
metric to the future horizon is the degenerate tensor
$r_h^2\sigma$ pulled back from its spherical space of generators.
For example, away from the bifurcation sphere the ingoing form is
\[
 \mathbf G=-F(\mathfrak r)\,dv^2
       +2\,dv\,d\mathfrak r+\mathfrak r^2\sigma,
 \qquad F(\mathfrak r)=1-\frac{2M}{\mathfrak r^p},
\]
and both of its first two terms restrict to zero at
$\mathfrak r=r_h$.  Kruskal coordinates give the same restriction at
the bifurcation sphere.  A tangent vector in the kernel of projection
onto the generator sphere would be null, so spacelikeness makes that
projection a local diffeomorphism.  Fullness makes it bijective, and
hence it is a diffeomorphism.
Pulling back the displayed degenerate tensor to the section therefore
gives a metric isometric to $r_h^2\sigma$, regardless of the location
of the section along individual generators.  This proves
Equation~\eqref{q:converse:eq:horizon-area}.

The null second fundamental form along the future horizon generator
vanishes as well: the tensor $r_h^2\sigma$ is constant along each
generator.  For a spacelike hypersurface ending on the section from
the exterior, the future outgoing null normal $\mathbf n+\nu$ is
a positive multiple of this generator, where $\nu$ points into the
hypersurface.  Scaling the generator multiplies its null second
fundamental form by the same positive factor.  Its trace is therefore
zero.  By the chosen convention that trace is
$H+\tr_S K=\theta_+$.
The same argument uses a regular future generator in Kruskal
coordinates at the bifurcation sphere, where the Killing field itself
vanishes.
\end{proof}

Proposition~\ref{q:converse:prop:adm} and
Lemma~\ref{q:converse:lem:horizon-area} prove the converse in
Theorem~\ref{thm:converse}.  Indeed, the assumed full-cut minimality
gives $A_{\min}(S)=A$, while the original charges give
\[
 \sqrt{E^2-|P|^2}=M
 =\frac12\left(\frac{A}{\omega}\right)^{p/k}.
\]
The argument did not impose a preferred original foliation or an
asymptotic graph hypothesis.  It derived the graph only on the tail
needed for the flux calculation; the compact part of the prescribed
embedding was arbitrary.

\subsection{Static and non-time-symmetric sharp examples}

The converse assumes full-cut minimality and outermostness.  To show
that its class is nonempty, we now construct slices for which these
properties can be checked for every enclosure.  A closed angular
form controls all cut areas, while the outward expansion of round
spheres excludes a wholly interior weakly trapped enclosure.

\begin{proposition}[Examples satisfying all full-cut hypotheses]
\label{q:converse:prop:examples}
For every $n\ge3$ and every $M>0$, the static exterior slice satisfies
all hypotheses of the rigidity subclass and realizes equality.
There are also smooth examples with $K\not\equiv0$, obtained by
nonconstant radial static-time graphs supported in any prescribed
annulus strictly outside the horizon.  These examples are vacuum,
complete with their boundary included, have the same horizon and
asymptotic end as the static slice, and satisfy
\[
 \Area_g(\Gamma)\ge\Area_g(S)
 \quad\hbox{for every full enclosing cut }\Gamma.
\]
They have no smooth compact weakly future-trapped enclosing
hypersurface lying wholly in the interior.
\end{proposition}

\begin{proof}
Fix $r_h<r_1<r_2$ and take a smooth function
$\chi\in C_c^\infty((r_1,r_2))$ with a nondegenerate interior critical
point.  Put $b=f(\mathfrak r)=\varepsilon\chi(\mathfrak r)$ and choose
$\varepsilon$ so that
\begin{equation}
 \sup_{\mathfrak r>r_h}|F(\mathfrak r)f'(\mathfrak r)|<1.
 \label{q:converse:eq:radial-spacelike}
\end{equation}
The zero choice gives the static slice.  The induced metric of any of
these graphs is
\begin{equation}
 g_f=A_f(\mathfrak r)\,d\mathfrak r^2+\mathfrak r^2\sigma,
 \qquad
 A_f=F^{-1}-F(f')^2
     =F^{-1}\bigl(1-F^2(f')^2\bigr)>0.
 \label{q:converse:eq:radial-metric}
\end{equation}
Thus the graph is spacelike.  Close to the horizon it is exactly the
static slice.  To check regularity there without relying on the
singular areal coordinate, use static proper distance $a\ge0$, for
which $d\mathfrak r/da=\sqrt F$.  At $a=0$,
\[
 \mathfrak r(a)
 =r_h+\frac{F'(r_h)}4a^2+O(a^4),\qquad F'(r_h)=\frac p{r_h}>0,
\]
and the metric is $da^2+\mathfrak r(a)^2\sigma$.
In Kruskal coordinates the static slice and its future normal extend
smoothly to the bifurcation sphere.  To see this, let the tortoise
coordinate $r_*$ satisfy $dr_*/d\mathfrak r=F^{-1}$.
The two null coordinates on the slice are opposite constant multiples of
$e^{\kappa r_*}$, where $\kappa=F'(r_h)/2$ and
$e^{\kappa r_*}=a$ times a smooth positive even function of $a$.
Their derivatives with respect to $a$ are nonzero at $a=0$.
Our perturbation is supported away from this collar, so its embedding
and normal have the same regularity.  The boundary is connected,
$K=0$ there, and its mean curvature is
$k\mathfrak r'(0)/r_h=0$.

The graph is unchanged at infinity as well.  Radial length to infinity
is infinite, and every bounded radial region is compact with the
boundary included; hence $(\Omega,g_f)$ is complete as a metric space
with that boundary.  The induced data have
$g_f-\delta=O_j(s^{-p})$ for every fixed $j$ in the static isotropic
end and $K=0$ there.  They therefore satisfy all the stated finite
derivative decay bounds for any $p/2<q<p$.

The ambient metric is vacuum in every required dimension.  This can
also be checked directly from its areal form: the only potentially
nonzero Ricci components are multiples of
\[
 F''+\frac{k}{\mathfrak r}F',\qquad
 (k-1)(1-F)-\mathfrak rF',
\]
and both vanish for $F=1-2M\mathfrak r^{-p}$ and $p=k-1$.
Gauss--Codazzi therefore gives $\mu=J=0$ on each graph, so the dominant
energy and integrability assumptions hold.  Its end charges are
$(E,P)=(M,0)$ by Lemma~\ref{q:converse:lem:boost} with $v=0$.
To verify non-time-symmetry when $\varepsilon\ne0$, take the chosen
critical point $\mathfrak r_*$ of $f$.  There $f'=0$ and the future
normal is $F^{-1/2}\partial_b$.  The defining second-form formula
then gives
\[
 K_{\mathfrak r\mathfrak r}(\mathfrak r_*)
   =\sqrt{F(\mathfrak r_*)}\,f''(\mathfrak r_*)\ne0.
\]

We prove the required area inequality for every full cut, including
disconnected cuts and all portions coinciding with the boundary.
Let $\pi_\sigma:\Omega\to S^k$ denote angular projection and set
\[
 \zeta=r_h^k\pi_\sigma^*(dA_\sigma).
\]
This is a smooth closed $k$-form on the entire exterior with boundary.
Its comass in the metric
Equation~\eqref{q:converse:eq:radial-metric}, meaning its maximum absolute
value on an orthonormal $k$-frame, is
$(r_h/\mathfrak r)^k\le1$.  Indeed it vanishes on the radial direction
and has that value on an oriented orthonormal angular frame.

Let $\Gamma=\partial D$ be a full cut as in
Definition~\ref{def:fullcuts}.  Choose $R$ beyond its entire boundary
and far enough that $D$ contains all $\mathfrak r\ge R$.
The intrinsic compact manifold
$D_R=D\cap\{\mathfrak r\le R\}$ has boundary exactly
$\Gamma\sqcup\{\mathfrak r=R\}$.  The second boundary is in the
interior of $D$ before truncation, so this operation introduces no
corner.  Orient $\Gamma$ by the normal pointing into $D$, and hence
toward its end.  Its orientation as a boundary of $D_R$ is the
opposite one.  Stokes' theorem gives
\begin{equation}
 \int_\Gamma\zeta
 =\int_{\{\mathfrak r=R\}}\zeta
 =\omega r_h^k.
 \label{q:converse:eq:full-cut-calibration}
\end{equation}
The entire intrinsic boundary appears in this formula, even when a
part of it lies on $S$.  Applying the comass bound to every component
and then summing gives
\[
 \Area_{g_f}(\Gamma)
 \ge\left|\int_\Gamma\zeta\right|
 =\omega r_h^k=\Area_{g_f}(S).
\]
Since $S$ itself is an admissible full cut, this proves both
$A_{\min}(S)>0$ and $A_{\min}(S)=\Area_{g_f}(S)$.

It remains to check the trapping exclusion.  Write
$d_f=1-F^2(f')^2>0$.  The future normal of the graph and the normal
to an outward round sphere within it are respectively
\[
 \mathbf n=\frac{F^{-1/2}}{\sqrt{d_f}}
             \bigl(\partial_b+F^2f'\partial_{\mathfrak r}\bigr),
 \qquad
 \nu=\frac{\sqrt F}{\sqrt{d_f}}
             \bigl(\partial_{\mathfrak r}+f'\partial_b\bigr).
\]
The round sphere has $H=k\nu(\mathfrak r)/\mathfrak r$ and
$\tr_{S_{\mathfrak r}}K=k\mathbf n(\mathfrak r)/\mathfrak r$.
Consequently
\begin{equation}
 \theta_+(S_{\mathfrak r})
 =\frac{k\sqrt F}{\mathfrak r\sqrt{d_f}}(1+Ff')>0
 \qquad(\mathfrak r>r_h).
 \label{q:converse:eq:round-expansion}
\end{equation}
The positivity follows from
Equation~\eqref{q:converse:eq:radial-spacelike}; it does not require a
smallness estimate for second derivatives of $f$.

Suppose a smooth compact embedded two-sided enclosing hypersurface
$\Sigma$ in the interior had $\theta_+\le0$ everywhere toward the
end.  At a maximum of $\mathfrak r$ on $\Sigma$, say
$\mathfrak r=\mathfrak r_0>r_h$, it is tangent to the round sphere
$S_{\mathfrak r_0}$ and lies on its inner side.  The ray with
$\mathfrak r>\mathfrak r_0$ reaches infinity without meeting
$\Sigma$, so the normal toward the end at this point is the positive
radial normal, even if $\Sigma$ is disconnected.  At the contact
point, tracing the restriction Hessian identity gives
\[
 0\ge\Delta_\Sigma\mathfrak r
  =\tr_{T\Sigma}\Hess_g\mathfrak r
            -H_\Sigma\,\nu(\mathfrak r).
\]
For the tangent round sphere the right-hand side is zero with
$H_\Sigma$ replaced by $H_{S_{\mathfrak r_0}}$.
Since $\nu(\mathfrak r)>0$, this implies
$H_\Sigma\ge H_{S_{\mathfrak r_0}}$.
The tangent spaces coincide, so their tangential traces of $K$ are
equal.  Equation~\eqref{q:converse:eq:round-expansion} now gives
$\theta_+(\Sigma)>0$ at the contact point, a contradiction.
All rigidity-subclass hypotheses have been checked.
\end{proof}

The examples exist for every $M>0$, and their common relation is
$A=\omega(2M)^{k/p}$.  Thus the coefficient $1/2$ is attained in every
dimension by both static and non-time-symmetric data.  Varying $M$
also fixes the exponent: an inequality
$M\ge c(A/\omega)^\beta$ with a universal $c>0$ for this entire scaled
family requires $\beta=p/k$, by considering both $M\downarrow0$ and
$M\to\infty$.
